\documentclass[10pt,twocolumn]{article}
\usepackage[margin=0.55in,top=0.45in,bottom=0.4in]{geometry}
\usepackage{graphicx}
\usepackage{booktabs}
\usepackage{enumitem}
\usepackage{xcolor}
\usepackage{titlesec}
\usepackage[hidelinks]{hyperref}

\setlength{\parskip}{3pt}
\setlength{\parindent}{0pt}
\setlist{nosep,leftmargin=*}

\titlespacing*{\section}{0pt}{6pt}{2pt}
\titlespacing*{\subsection}{0pt}{4pt}{1pt}
\titleformat{\section}{\normalfont\bfseries\normalsize}{\thesection}{0.5em}{}
\titleformat{\subsection}{\normalfont\bfseries\small}{\thesubsection}{0.5em}{}

\begin{document}

\twocolumn[{
\begin{center}
{\Large\bfseries Autoreason: Autoresearch for Subjective Domains}\\[4pt]
{\normalsize Nous Research}\\[2pt]
{\small\url{https://github.com/NousResearch/autoreason}}\\[8pt]
\end{center}
}]

\section{Problem}

LLMs used for iterative refinement exhibit compounding failure modes: \textbf{sycophancy} (strengthens whatever you hand it), \textbf{overcriticism} (always finds problems even when work is sound), \textbf{overcompromise} (hedges everything when synthesizing), \textbf{authorship bias} (defends its own prior choices), \textbf{scope drift} (adds complexity each iteration), and \textbf{context collapse} (output diverges from original intent with no reversion mechanism). The output is shaped more by how you prompt than by what is actually better.

\section{Method}

Autoreason extends autoresearch to subjective domains by constructing a fitness function from independent blind evaluation---the same way science uses peer review where math can use proofs.

Every role is a fresh, isolated agent with no shared context. The loop:

\begin{enumerate}
\item \textbf{Version A} --- the incumbent (unchanged)
\item \textbf{Strawman} --- fresh agent attacks A, finds problems only
\item \textbf{Version B} --- fresh agent revises A using valid criticisms
\item \textbf{Version AB} --- fresh agent synthesizes best of A and B
\item \textbf{Judge panel} --- 3 blind judges rank A, B, AB against the original task prompt (randomized labels, Borda count, incumbent wins ties)
\end{enumerate}

Winner becomes new A. Loop repeats until A survives 2 consecutive passes.

\textbf{A} represents conservatism, \textbf{B} represents adversarial editing, \textbf{AB} approximates objectivity. The judge panel decides which framing produced the best result, with no knowledge of which is which.

\section{Results}

\subsection{5-Way Baseline Comparison}

All methods start from the same initial output (847 words), 15 passes, same model (Claude Sonnet 4). 7-judge blind panel, randomized labels, Borda count.

\begin{table}[h]
\centering
\small
\begin{tabular}{@{}lccr@{}}
\toprule
\textbf{Method} & \textbf{Borda} & \textbf{1st} & \textbf{Words} \\
\midrule
\textbf{Autoreason} & \textbf{35/35} & \textbf{7/7} & 1,800 \\
Conservative & 21 & 0 & 862 \\
``Improve this'' & 18 & 0 & 2,116 \\
Harsh critic & 18 & 0 & 1,961 \\
Critique \& revise & 13 & 0 & 2,507 \\
\bottomrule
\end{tabular}
\end{table}

Autoreason won unanimously. Conservative (barely changed the initial output) beat all three iterative baselines---doing almost nothing is better than iterating without structural protection. Critique-and-revise, the standard approach, came last.

\subsection{Qualitative Improvement}

The initial output is a generic startup playbook with round-number targets and no validation. After 14 passes, the autoreason output contains quantified pain (\$15K/incident $\times$ 6/year), team-based pricing matching buying motions (\$1,499/mo vs \$49/user), validation from 50+ interviews with 75\% pilot success, competitive positioning against named tools, and unit economics (CAC \$2K, LTV \$54K). Unrealistic revenue targets were killed (\$100K $\to$ \$25K MRR).

\subsection{Convergence}

The loop reached 2 consecutive incumbent wins at pass 15. Pass 15 vs pass 25 (second convergence): 6--1 for pass 15. First convergence is the quality ceiling.

\subsection{Key Findings}

\begin{itemize}
\item \textbf{Fresh agents prevent authorship bias.} Author B re-emerged as winner at passes 17--21 after 15 passes of irrelevance. A persistent agent would have learned to defer.
\item \textbf{Bloat/prune oscillation.} On ambiguous tasks, AB adds complexity while B prunes it. Word counts: 847 $\to$ 1800 $\to$ 1644 $\to$ 1758. Signals the task is underdetermined along the scope dimension.
\item \textbf{Conservative tiebreak is load-bearing.} Incumbent wins ties. Removed 3 unnecessary churn events in 26 passes.
\item \textbf{Judges need the initial output as baseline.} Without it, autoreason vs adversarial was 3--2. With it, 7--0. Reference points expose drift.
\end{itemize}

\section{Related Work}

\textbf{Autoresearch} (Karpathy) uses val\_bpb as an objective fitness function for autonomous experiments. Autoreason extends this to domains without objective metrics. \textbf{SlopCodeBench} (Orlanski et al., 2026) shows agents produce progressively worse code each iteration---prompting doesn't fix it. \textbf{ACE} (Zhang et al., ICLR 2026) documents context collapse in iterative rewriting. \textbf{LLM Council} (Karpathy) uses multi-model panels for single-turn evaluation; autoreason adds iteration with role separation.

\section{Open Questions}

Does autoreason improve outcomes where metrics exist but are insufficient (code that passes tests but degrades)? Do mixed-model judge panels converge faster? Does the loop reach the same optimum across independent runs?

\vfill
{\small Code, data, and artifacts: \url{https://github.com/NousResearch/autoreason}}

\end{document}
